En este trabajo se presenta un estudio comparativo del rendimiento
practico frente a la complejidad teórica de algoritmos de
ordenamiento de arreglos unidimensionales (MergeSort, QuickSort,
PatienceSort y Sort) y de multiplicación de matrices cuadradas
(Naive y Strassen). Las pruebas fueron implementadas en C++ y
configuradas para medir el tiempo de ejecución y uso de memoria
auxiliar (\textit{Peak RSS}) de cada algoritmo bajo diversas condiciones de entrada.
Los resultados indican que Sort es la alternativa más
eficiente e in-place para ordenamiento, mientras que QuickSort muestra
una alta sensibilidad a recibir de entrada arreglos ordenados, aumentando
su tiempo de ejecución a $O(n^2)$ por la partición de Lomuto.
Para multiplicación de matrices en dimensiones $n \le 1024$,
el algoritmo Naive superó a Strassen tanto en tiempo como en
memoria ($1300\text{ KB}$ frente a $6700\text{ KB}$),
evidenciando que el \textit{overhead} recursivo predomina
sobre la ventaja asintótica en dimensiones pequeñas y medianas.